\section{Problem Formulation}

D'Angelo et al.~\cite{dangelo2025forsid} formalize the Forget-Set Identification (ForSId) problem: given a training set $D$, a model $M_D$, an unwanted set $D_u$, and a wanted set $D_w$, find a forget set $D_f \subseteq D$ whose removal maximally preserves model behavior on $D_w$ while altering it on $D_u$. ForSId requires model access to compute per-sample influence via training-data gradients. More fundamentally, its input is \emph{behavioral} evidence---samples whose predictions should change---not metadata-level queries such as ``forget all data by author~$a_i$.''

When author $a_i$ initiates a withdrawal request, the target forget set is $F = \{l \in D \mid \text{author}(l) = a_i\}$. Computing $F$ precisely without model access is the challenge ob addresses. Dataset-level provenance can only answer ``data from author~$a_i$ exists in~$D$,'' forcing two extremes: delete the entire dataset or take no action. Line-level provenance resolves this by answering ``author~$a_i$ contributed $n$ lines,'' enabling targeted deletion---as we show in Section~\ref{sec:revocation-precision}, record-level revocation reduces over-deletion by up to 101$\times$ compared to dataset-level approaches.

From this gap, we derive three design requirements. \emph{Precision}: locate specific data lines, not files or datasets. \emph{Verifiability}: provenance must be independently auditable, not merely stored internally. \emph{Minimal invasiveness}: integration must require only a few lines of code---DLProv~\cite{pina2025dlprov} serves as a cautionary example, requiring extensive script instrumentation despite low runtime overhead.
